% ID: SEL-D0064-OB-E04
% ESTADO: validada. S064.
\textbf{a)} Al $98\%$ queda un $1\%$ en cada cola. Con desviación conocida, el margen es $z_{0{,}99}\,11/\sqrt n$. Como debe ser estrictamente menor que $1$ cm,
\[
n>(11z_{0{,}99})^2\approx654{,}84.
\]
Con el cuantil sin redondear a dos decimales, el mínimo es $\boxed{655\text{ alumnos}}$.

\textbf{b)} Para muestreo aleatorio simple sin reposición, las muestras equiprobables $\{10,12\}$, $\{10,17\}$ y $\{12,17\}$ tienen medias $11$, $13{,}5$ y $14{,}5$. Su media es $\boxed{13}$ y la desviación típica es
\[
\sqrt{\frac{(11-13)^2+(13{,}5-13)^2+(14{,}5-13)^2}{3}}
=\boxed{\sqrt{13/6}}\approx1{,}472.
\]
Si se usan dos extracciones independientes con reposición, hay nueve pares ordenados equiprobables. Sus medias son $10,11,13{,}5,11,12,14{,}5,13{,}5,14{,}5,17$. El centro sigue siendo $13$ y la desviación pasa a $\sqrt{13/3}$, al dividir la varianza poblacional $26/3$ entre dos.
